\documentclass[12pt,a4paper]{article}
\usepackage[T1]{fontenc}
\usepackage[utf8]{inputenc}
\usepackage[english]{babel}
\usepackage{enumitem}
\usepackage{geometry}
\usepackage{hyperref}
\geometry{margin=2.5cm}

\title{Evaluation Chapter: Validation of n8n Workflows}
\author{}
\date{}

\begin{document}

\maketitle

\section{Objective}

The goal of the evaluation is to systematically compare different validation approaches for n8n workflows. The study investigates which methods reliably detect faulty workflows, from which level of complexity errors begin to be overlooked, and whether non-LLM-based validation procedures can also be used for LLM tasks within a workflow.

\section{Core Idea}

Each test case consists of:

\begin{itemize}
  \item a natural-language business description,
  \item a correct reference workflow,
  \item several intentionally faulty variants,
  \item one or more validators.
\end{itemize}

For each validator, the evaluation measures whether correct and incorrect workflows are classified correctly and which error types are detected or overlooked.

\section{Dataset Design}

A realistic dataset can consist of \texttt{100} base workflows.

Recommended distribution:

\begin{itemize}
  \item \texttt{40} classical sequential business workflows,
  \item \texttt{30} workflows with branching and parallelism,
  \item \texttt{20} more complex exception and approval processes,
  \item \texttt{10} workflows with \texttt{1--2} LLM nodes.
\end{itemize}

For each base workflow, several variants are generated:

\begin{itemize}
  \item \texttt{1} correct workflow,
  \item \texttt{1} structural error variant,
  \item \texttt{1} control-flow error variant,
  \item \texttt{1} semantic error variant,
  \item for LLM workflows, additionally \texttt{1} LLM error variant.
\end{itemize}

This results in approximately \texttt{410} to \texttt{500} test objects.

\section{Error Classes}

\subsection{Structural Errors}

\begin{itemize}
  \item missing node,
  \item wrong connection,
  \item dead path,
  \item unreachable node,
  \item missing merge.
\end{itemize}

\subsection{Control-Flow Errors}

\begin{itemize}
  \item swapped IF branch,
  \item missing parallelism,
  \item wrong order,
  \item incorrectly placed wait or retry step.
\end{itemize}

\subsection{Semantic Errors}

\begin{itemize}
  \item business-wise wrong step,
  \item missing required step despite formally valid workflow,
  \item wrong message recipient,
  \item invoice booked before approval.
\end{itemize}

\subsection{LLM-Related Errors}

\begin{itemize}
  \item LLM node at the wrong position,
  \item LLM node with the wrong purpose,
  \item LLM output is accepted without further validation,
  \item LLM makes a decision that should be checked deterministically.
\end{itemize}

\section{Validators}

\subsection{Static Validator}

Checks:

\begin{itemize}
  \item JSON structure,
  \item node types,
  \item connections,
  \item reachability,
  \item dead paths.
\end{itemize}

Strengths:

\begin{itemize}
  \item strong for technical and structural errors.
\end{itemize}

Weaknesses:

\begin{itemize}
  \item weak for business semantics.
\end{itemize}

\subsection{Rule-Based Business Validator}

Checks:

\begin{itemize}
  \item required steps,
  \item order,
  \item branches,
  \item parallelism,
  \item target-vs-actual comparison with a reference specification.
\end{itemize}

Strengths:

\begin{itemize}
  \item strong for modellable business logic.
\end{itemize}

Weaknesses:

\begin{itemize}
  \item limited for open meanings and language variation.
\end{itemize}

\subsection{LLM-as-a-Judge}

Input:

\begin{itemize}
  \item business description,
  \item workflow JSON or an abstracted workflow view.
\end{itemize}

Output:

\begin{itemize}
  \item \texttt{valid} or \texttt{invalid},
  \item explanation,
  \item optionally error category and confidence.
\end{itemize}

Strengths:

\begin{itemize}
  \item better for semantic contradictions.
\end{itemize}

Weaknesses:

\begin{itemize}
  \item potentially unstable for subtle structural details,
  \item more expensive.
\end{itemize}

\subsection{Hybrid}

Combines:

\begin{itemize}
  \item static and rule-based pre-validation,
  \item LLM-based residual validation for semantic deviations.
\end{itemize}

Strengths:

\begin{itemize}
  \item likely the most robust overall.
\end{itemize}

Weaknesses:

\begin{itemize}
  \item higher complexity,
  \item higher cost.
\end{itemize}

\section{Additional Research Question: Complexity}

In addition to overall performance, the evaluation investigates whether validation quality decreases with increasing workflow complexity.

Workflow complexity is operationalized through:

\begin{itemize}
  \item number of nodes,
  \item number of edges,
  \item number of branches,
  \item number of merges,
  \item control-flow depth,
  \item number of external integrations,
  \item number of LLM nodes.
\end{itemize}

Analysis:

\begin{itemize}
  \item performance per complexity level,
  \item errors that are only overlooked from complexity level \texttt{x} onward,
  \item differences between classical workflows and workflows containing LLM nodes.
\end{itemize}

\section{Additional Research Question: LLM Tasks in the Workflow}

A separate part of the evaluation examines LLM tasks within a workflow.

The aim is not to validate the full truth of the LLM output. Instead, the following aspects are checked:

\begin{itemize}
  \item whether the LLM node is placed at the correct business position,
  \item whether its purpose matches the description,
  \item whether its output is validated further,
  \item whether critical decisions are not based blindly on the LLM,
  \item whether the output flows into the correct branch.
\end{itemize}

Two approaches are compared:

\begin{itemize}
  \item \texttt{LLM-as-a-Judge} for the entire LLM part,
  \item a \texttt{rule-based approximation} without a judge model.
\end{itemize}

The rule-based approximation can build on techniques such as those in \texttt{validate\_n8n\_example.ps1}:

\begin{itemize}
  \item keyword matching,
  \item stopword filtering,
  \item start, end, and branch detection,
  \item checking expected node names,
  \item comparison of parallelism and alternatives.
\end{itemize}

This makes it possible to investigate whether, for certain LLM tasks, a relatively simple and explainable validation is already sufficient and where a real judge model remains necessary.

\section{Handling External Integrations}

An important part of the test setup concerns external systems such as:

\begin{itemize}
  \item email servers,
  \item HTTP APIs,
  \item databases,
  \item ERP or CRM systems.
\end{itemize}

For workflow correctness validation, it is generally not sensible to run these integrations by default against real production or live systems.

\subsection{Principle}

For the main evaluation, external integrations should preferably be \texttt{mocked} or \texttt{simulated}.

\subsection{Reasoning}

\begin{itemize}
  \item Tests become reproducible and controllable.
  \item Side effects such as real emails or real bookings are avoided.
  \item Authentication and network problems do not distort the validation.
  \item Errors can be attributed more clearly to the workflow instead of a third-party system.
  \item Privacy and security risks are reduced.
\end{itemize}

\subsection{Specifically for Email Servers}

For email nodes, mocking or simulation is the most sensible choice for most test cases.

Instead of real delivery, the evaluation can check:

\begin{itemize}
  \item whether an email would have been sent,
  \item to which recipient,
  \item with which subject,
  \item with which expected business content,
  \item under which workflow condition.
\end{itemize}

Possible implementations:

\begin{itemize}
  \item fake SMTP server,
  \item local mail catcher,
  \item test endpoint that logs sending events,
  \item replacement of the email node by a logging or storage node in a test variant.
\end{itemize}

\subsection{When Real Sending Can Still Be Useful}

A small integration sample with real or near-real test infrastructure can additionally be useful to demonstrate external connectivity.

However, this should only play a minor role and should not dominate the main evaluation.

\subsection{Methodological Decision}

The recommended setup is:

\begin{itemize}
  \item main evaluation with mocked or simulated integrations,
  \item optional small integration validation with real test systems.
\end{itemize}

This ensures that not only structure but also the intended behavior of integration nodes can be validated without compromising scientific comparability through external instability.

\section{Test Procedure}

For each test object:

\begin{enumerate}
  \item load the reference description and reference specification,
  \item apply the validator to the workflow,
  \item store the output:
  \begin{itemize}
    \item \texttt{valid} or \texttt{invalid},
    \item detected errors,
    \item confidence, if available,
    \item runtime,
    \item token cost, if an LLM is used,
  \end{itemize}
  \item compare the result with the ground truth.
\end{enumerate}

\section{Metrics}

Mandatory metrics:

\begin{itemize}
  \item Accuracy,
  \item Precision,
  \item Recall,
  \item F1,
  \item False Positive Rate,
  \item False Negative Rate.
\end{itemize}

Additional analyses:

\begin{itemize}
  \item per error class,
  \item per complexity level,
  \item for workflows with and without LLM nodes,
  \item average runtime,
  \item average token cost.
\end{itemize}

\section{Possible Result Presentations}

To ensure that the evaluation does not remain purely textual, the results should be prepared using several complementary representations.

\subsection{Tables}

\begin{itemize}
  \item table with overall values per validator: Accuracy, Precision, Recall, F1,
  \item table per error class: structural, control-flow, semantic, LLM-related,
  \item table per complexity level: low, medium, high,
  \item table for workflows with and without LLM nodes,
  \item table for average runtime and average token cost.
\end{itemize}

\subsection{Bar Charts}

\begin{itemize}
  \item bar chart: \texttt{F1} per validator,
  \item bar chart: \texttt{Recall} per error class and validator,
  \item bar chart: \texttt{False Negative Rate} for complex workflows,
  \item bar chart: comparison of LLM workflows and classical workflows.
\end{itemize}

\subsection{Line Charts}

\begin{itemize}
  \item line: validator performance over increasing node count,
  \item line: performance over increasing branch count,
  \item line: recall as a function of workflow depth,
  \item line: error rate from complexity level \texttt{x}.
\end{itemize}

\subsection{Heatmaps}

\begin{itemize}
  \item heatmap: validators by error classes,
  \item heatmap: validators by complexity levels,
  \item heatmap: LLM node types by detection rate.
\end{itemize}

\subsection{Boxplots}

\begin{itemize}
  \item boxplot: runtime per validator,
  \item boxplot: token costs per LLM-based validator,
  \item boxplot: variation in performance across different domains.
\end{itemize}

\subsection{Confusion Matrices}

\begin{itemize}
  \item confusion matrix per validator for \texttt{valid} vs. \texttt{invalid},
  \item separate matrices for the full set and the LLM subset.
\end{itemize}

\subsection{Case-Based Visualizations}

\begin{itemize}
  \item small workflow figures for selected good and bad examples,
  \item plus a table with:
  \begin{itemize}
    \item ground truth,
    \item validator decision,
    \item detected errors,
    \item overlooked errors.
  \end{itemize}
\end{itemize}

\section{Possible Result Lists in the Results Chapter}

In the results chapter, findings could be presented roughly in the following order:

\begin{enumerate}
  \item Overall comparison of all validators
  \begin{itemize}
    \item Which validator achieves the best overall performance?
    \item Which differences are statistically or practically relevant?
  \end{itemize}
  \item Results by error class
  \begin{itemize}
    \item Who detects structural errors best?
    \item Who detects control-flow errors best?
    \item Who detects semantic errors best?
    \item How do validators perform on LLM-related errors?
  \end{itemize}
  \item Results by complexity
  \begin{itemize}
    \item From which node count does performance decrease?
    \item From which branching depth are errors frequently overlooked?
    \item Which method is most robust for complex workflows?
  \end{itemize}
  \item Results for LLM workflows
  \begin{itemize}
    \item How strongly does validation quality decrease for LLM nodes?
    \item Where is the rule-based approximation sufficient?
    \item Where is a judge model clearly required?
  \end{itemize}
  \item Cost-benefit analysis
  \begin{itemize}
    \item How much does the hybrid approach improve performance?
    \item Does the performance gain justify the additional token cost?
    \item Which method is most efficient for practical use?
  \end{itemize}
  \item Qualitative error analysis
  \begin{itemize}
    \item typical missed errors per validator,
    \item typical false alarms per validator,
    \item concrete examples of edge cases.
  \end{itemize}
\end{enumerate}

\section{Expected Non-Trivial Results}

It is not expected that a single validator will dominate in all dimensions.

Likely patterns:

\begin{itemize}
  \item the static validator performs best for structural errors,
  \item the rule-based validator performs well for clear invariants,
  \item \texttt{LLM-as-a-Judge} performs better for semantic deviations,
  \item the hybrid approach is overall the most robust but more expensive,
  \item for LLM tasks, the reliability of all methods decreases, but not equally strongly.
\end{itemize}

These differentiated differences are what make the evaluation scientifically interesting.

\section{Conclusion of the Test Setup}

This setup enables:

\begin{itemize}
  \item a fair comparison of multiple validation approaches,
  \item an investigation of error detection as a function of complexity,
  \item a targeted analysis of whether LLM tasks in workflows are partially validatable without a judge model.
\end{itemize}

This is likely to produce heterogeneous results instead of a trivial statement such as \texttt{everything is validatable} or \texttt{nothing is validatable}.

\end{document}
